% =====================================================================
%  Harald Høffding, „Nye Undersøgelser om den primitive Familje“, review
%  of C. N. Starcke, Die primitive Familie in ihrer Entstehung und
%  Entwickelung (Leipzig: Brockhaus, 1888).
%  Tilskueren. Maanedsskrift, 5. Aargang (1888), pp. 787--795.
%  TRANSCRIPTION (Danish) -- diplomatic, from the page images.
%
%  ---- SOURCE -----------------------------------------------------------
%  Internet Archive item tilskuerenmaaned1888nngoog (Google Books scan of
%  the whole volume, Tilskuerenmaaned1888nngoog.pdf, 978 PDF pp.,
%  30586820 bytes, sha256
%  d20291f6efa41ced8785c2b0a9f72a32b9152935e1807db08523e519dc6f44ce; no
%  embedded text layer), downloaded 2026-09-27:
%  https://archive.org/details/tilskuerenmaaned1888nngoog
%  The article is volume PDF pp. 799--807.  Kept as an extract (qpdf) in
%  ~/bibliotek/Høffding, Harald/:
%    1888-primitive-familie-tilskueren-5.pdf  sha256 6d6a8628...
%  Each page one 600 ppi bitonal image.  Antiqua (roman).
%
%  ---- PAGE MAP ---------------------------------------------------------
%  Volume PDF = printed + 12; verified by eye on the folios.
%
%  ---- CONVENTIONS ------------------------------------------------------
%  Diplomatic; 1888 orthography as printed.  Letterspacing -> \emph.
%  Quotation marks as printed.  Footnotes *) at the mark.  Printer's
%  errors transcribed as printed and logged in a % comment at the site.
%
%  ---- LAYOUT -----------------------------------------------------------
%  p. 787 opens a fresh page (no folio; „53*“ is the sheet signature):
%  the bold head, the reviewed book's title in small type (the author's
%  name letterspaced), a short rule, and the text under an enlarged
%  initial.  No numbered sections, no footnotes.  Quotations in
%  guillemets »...« (8 pairs, balanced).  Emphasis (letterspacing,
%  measured within the line): 16 runs.  Closes (p. 795) with the
%  signature set right in small caps and a double rule.
%
%  Printer's errors logged at the site: p. 788 „nøgterne“ (for
%  nøgternt); p. 789 „Untermehung“ (for Untersuchung, in a German
%  title) and „Promiskniteten“ (elsewhere Promiskuiteten); p. 792
%  „Kønsinktinktet“; p. 795 a full stop before lower-case „skilles“.
%  The three misspellings confirmed at the image in the caller.  Broken
%  sorts noted (p. 788 „derved“, p. 791 „Hustru“, p. 795 „endnu“).
%
%  ---- ACCURACY  [2026-09-27] ------------------------------------------
%  The scan has no text layer.  Word-accuracy: the batch was diffed
%  (ocrdiff.py --frag --ocr) against an independent witness, the Internet
%  Archive's ABBYY OCR of the volume (Tilskuerenmaaned1888nngoog_djvu.txt),
%  cut to pp. 787-795 one page per form feed: 22 candidates, 0
%  transcription errors corrected, 22 the witness's fault, 0 unresolved.
%  The transcriber read every word at the image (300 dpi; 600 dpi for any
%  doubt), so the witness is independent of the words transcribed.
%  SECOND READING of all nine pages, word by word and mark by mark at the
%  image (300 dpi; 600 dpi for any doubt), by a reader who had transcribed
%  nothing and used no OCR: no discrepancy (0 word, 0 other).  The reader
%  settled the open point (p. 788: a comma after „Folkeslag“) and
%  confirmed the full stop after „Bog“ (p. 795), the 8 guillemet pairs,
%  the 16 letterspaced runs and the small-caps signature.
% =====================================================================
\documentclass[12pt,a4paper,openany]{book}
\usepackage[utf8]{inputenc}
\usepackage[T1]{fontenc}
\usepackage{textalpha}
\usepackage{libertinus}
\usepackage{libertinust1math}
\usepackage{graphicx}    % for \reflectbox (the old Danish ɔ: id-est mark)
\DeclareUnicodeCharacter{0254}{\reflectbox{c}}  % ɔ (U+0254) = reversed c
\usepackage[danish]{babel}
\usepackage[protrusion=true,expansion=true]{microtype}
\usepackage{setspace}
\usepackage[margin=1.2in]{geometry}
\usepackage{fancyhdr}
\setlength{\headheight}{28pt}
\addtolength{\topmargin}{-13.5pt}

% Footnotes as the printing sets them: *).
\renewcommand{\thefootnote}{*)}

% The article's numbered sections: the number alone, centred on its line.
\newcommand{\secnum}[1]{\par\begin{center}#1\end{center}\par\nobreak}

\usepackage{hyperref}
\hypersetup{
  pdftitle={Nye Undersøgelser om den primitive Familje},
  pdfauthor={Harald Høffding},
  hidelinks
}

\pagestyle{fancy}
\fancyhf{}
\fancyhead[LE,RO]{\thepage}
\fancyhead[RE]{\textit{Harald Høffding: Nye Undersøgelser om den primitive Familje}}
\fancyhead[LO]{\textit{Tilskueren, 5.\ Aargang (1888)}}

\setstretch{1.3}

\title{\textbf{Nye Undersøgelser om den primitive Familje}}
\author{Harald Høffding}
\date{Tilskueren. Maanedsskrift, 5.\ Aargang (1888), pp.~787--795}

\usepackage{marginnote}
\newcommand{\opage}[1]{\marginnote{\footnotesize\textit{[#1]}}}
\newcommand{\apage}[1]{\marginnote{\footnotesize\textit{[#1]}}}

\begin{document}
\maketitle
\thispagestyle{empty}
\clearpage

% --- p. 787 ---
% p. 787: the upper part of the page is blank (the article opens a new page; no running head, no folio -- the „53*“ at the foot is the sheet signature, not transcribed); the head „Nye Undersøgelser om den primitive Familje.“ in bold display antiqua on one line; below it the bibliographic line in small antiqua, centred, in three lines, „Starcke“ letterspaced (gaps 5, 6, 6, 5, 6, 7 px at 300 dpi against 1--3 px solid); a short centred rule (measured c. 1.3 cm, 152 px at 300 dpi); the text opens with an enlarged initial D, no indent
\opage{787}\begin{center}\textbf{\large Nye Undersøgelser om den primitive Familje.}\end{center}

\begin{center}\small C. N. \emph{Starcke}: Die primitive Familie in ihrer Entstehung und Entwickelung. Leipzig. F. A. Brockhaus. 1888. (Internationale wissenschaftliche Bibliothek. 66. Bind).\end{center}

\begin{center}\rule{1.3cm}{0.4pt}\end{center}

Det vil næppe undre nogen, at en filosofisk Forfatter giver
sig af med Undersøgelser om Familjens Historie. Naar Filosofi
skal være andet og mere end Konstruktioner ud fra den subjektive
Bevidsthed, og naar særlig Etik skal være noget mere end
Opstilling af Postulater, maa der skaffes et faktisk, og det vil for
Etikens Vedkommende sige: et historisk Materiale til Veje, som
paa mange Maader kan belyse Principernes Værdi og Mulighederne
for deres Anvendelse. Der maa ganske vist — i det mindste
efter min Opfattelse — skelnes bestemt mellem historisk og filosofisk
Etik. Hin kan kun fremstille det etiske Liv og de etiske
% emphasis: „faktisk“ (gaps 6, 7, 6, 9, 7 px at 300 dpi) and „Vurderingen“ (gaps 6, 7, 8, 7, 7, 8, 7, 9, 6, 6 px) letterspaced, against 1--4 px solid
Institutioner som de \emph{faktisk} have udviklet sig under visse Forhold,
men en derfra forskellig Opgave er \emph{Vurderingen} af dem,
Bedømmelsen af deres Betydning for den fortsatte, bevidste Udvikling
af det menneskelige Liv. Men den bevidste Indgriben
forudsætter stedse Udviklingsmuligheder og Spirer, som uvilkaarlig
ere blevne til. Kendskab til Fortidens historiske Livsformer er
derfor en nødvendig Betingelse for en frugtbar Behandling af
Etiken. Og det er ikke altid sagt, at Historikere og Etnologer
lægge Stoffet til Rette saaledes som Filosofien behøver det. Dertil
hører ikke blot historisk Kritik, men ogsaa filosofisk Blik paa de
afgørende Punkter. Ligesom Retshistorien ikke kan undvære
Medvirken fra Juristernes Kreds, saaledes vil filosofisk Dannelse
være en nødvendig Forudsætning for Behandlingen og Fremstillingen
af den historiske Etik.

Særlig Opmærksomhed har i de senere Aar den primitive
Familje, d. v. s. Familjen, saaledes som den fremtræder paa de
% --- p. 788 ---
\opage{788}laveste os bekendte menneskelige Udviklingstrin, vakt. Medens
man tidligere blev staaende ved den patriarkalske Familje som
man kendte den fra de gammeltestamentlige Fortællinger og fra
de indiske Digte og Lovbøger, er nu det betydningsfulde Skridt
sket, at den historiske eller sammenlignende Etik ogsaa begyndte
at tage Hensyn til de vilde Folkeslags Familjeforhold. Man saa
nu, at man ikke havde Ret til at betragte det hebræiske og
indiske Udviklingstrin som det mest primitive. Der kom her en
lignende Revolution i Tankegangen, som da Sprogvidenskaben gik
ud over Sanskrit og tog Hensyn til de vilde Nationers Sprog,
% as printed: „Henstillinger“ (the sense suggests „Forestillinger“); not corrected
eller da Mytologien inddrog de Vildes religiøse Henstillinger under
sine Betragtninger. Man stødte paa nye og overraskende Kendsgerninger
% broken letter: the r of „derved“ has lost its shoulder (reads „deıved“); transcribed as the word printed
og ledtes derved til dristige Hypoteser. Dr. Starckes
Skrift, som bygges paa et overordentlig grundigt Kendskab til det
foreliggende Materiale, og som bærer Vidnesbyrd om kritisk Ævne
og sundt psykologisk Blik, vil ganske sikkert komme til at indtage
en betydelig Plads i Literaturen om det interessante Æmne,
det behandler. Det forsmaar de ydre Tillokkelser, der for mange
kunde have ligget i de pikante Træk og underholdende Smaahistorier,
som Æmnet og Kilderne frembyde i rigelig Mængde.
% printer's error: „nøgterne“ (for „nøgternt“, beside „strængt“); transcribed as printed
Det er nøgterne og strængt i sin Form, men tiltrækkende og lærerigt
for den, der vil orienteres i Æmnet og have Vejledning til at
tage Stilling overfor de dristige, ofte vilde Paastande og Konsekvenser,
hvortil det har givet Anledning. — Lad mig forsøge
kortelig at angive, hvorom det drejer sig, og hvilket Standpunkt
Dr. Starcke indtager.

Man er bleven opmærksom paa, at hos en Mængde Folkeslag
rundt om paa Jorden regnes Slægtskabet ikke gennem Faderen,
men gennem Moderen, saaledes at en Mands nærmeste Arvinger
ikke ere hans Sønner, men hans Søsters Børn. I alle Verdensdele
findes denne Skik. Det laa da nær at antage, at der var en
fælles Aarsag til den. Den store Løshed og Frihed i kønslig
% p. 788: the mark after „Folkeslag“ at the distorted right margin is a comma (tail below the baseline at 600 dpi; settled by the second reading)
Henseende, som ligeledes findes hos de fleste vilde Folkeslag,
gjorde, at man kom til at tænke paa den gamle Sætning: pater
incertus, mater semper certa. (Man kan være uvis om, hvem
Faderen er, men Moderen er altid sikker). Man antog, at den
oprindelige Tilstand i kønslig Henseende havde været en fuldstændig
Ubundethed og en stadig Skiften, en Parring i Flæng
(Promiskuitet), som gjorde det umuligt at henføre Børnene til
bestemte Fædre og nødte til at lægge Kvindelinjen til Grund.
% --- p. 789 ---
\opage{789}Den, som først udtalte den Hypotese, at der havde været en Tid,
hvor Slægtskabet regnedes gennem Kvinden, en Ordning, der betegnede
et Fremskridt ved at afløse eller dog regulere Promiskuiteten,
% emphasis: „Bachofen“ letterspaced (gaps 6, 6, 8, 6, 2, 6 px at 300 dpi)
var \emph{Bachofen} i hans Værk »Das Mutterrecht. Eine
% printer's error: „Untermehung“ (for „Untersuchung“); transcribed as printed
Untermehung über die Gynaikokratie der alten Welt nach ihrer
religiøsen und rechtlichen Natur« (Stuttgart 1861). Han satte
Kvindelinjens Fremhersken i en vis Periode i Forbindelse med
en dristig mytologisk Teori og mente, at Kvinden i denne Periode
havde været højt anset, og at hendes Herredømme (Matriarkatet)
havde haft stor kulturhistorisk Betydning ved at mildne Sæderne
og hæmme de dyriske Instinkter. En Støtte for sin Teori mente
han at finde i Polyandriet, den Form for kønsligt Samliv, hvor
en Kvinde er flere Mænds Hustru. Han ansaa det for en nødvendig
Mellemform, gennem hvilken Menneskeslægten har arbejdet
% printer's error: „Promiskniteten“ (n for u; elsewhere „Promiskuiteten“); transcribed as printed
sig ud af Promiskniteten og nærmet sig Monogamiet. — Blandt
% emphasis: „Mc Lennan“ (gaps 7; 6, 6, 6, 7, 6 px) and „Lubbock“ (gaps 7, 8, 6, 6, 6, 7 px) letterspaced; the second „Mc Lennan“ below is solid
dem, der have opstillet lignende Teorier, maa \emph{Mc Lennan} og
\emph{Lubbock} nævnes. Disse Forskere gaa frem med større Kritik
og mindre Mystik end Bachofen, og især Mc Lennan søger paa
en skarpsindig Maade at vise, hvorledes mange Sæder og Skikke
paa højere Udviklingstrin, som ellers vilde være uforklarlige, kunne
forstaas som Rester (survivals) af Kvindelinje eller Polyandri med
Promiskuitet som Baggrund. Saadanne Skikke ere f. Eks. Leviratet,
den Skik, at en Mand var forpligtet til at ægte sin afdøde
Broders Enke, saaledes at Børnene, som udsprang af dette Ægteskab,
dog regnedes for den afdødes Børn, og Niyoga, den Skik,
at Ægtemanden (endnu hos Inder, Spartaner, ja hos Middelalderens
Germaner) kunde overlade sin Hustru til andre Mænd og dog
betragtede de saaledes avlede Børn som sine egne.

Efter denne Teori vilde Billedet af Menneskeslægtens Udvikling,
hvad Forholdet mellem Mand og Kvinde angaar, ganske
vist afrunde sig paa en klar og simpel Maade. Fra et oprindeligt
ubundet Fællesskab (Promiskuiteten) vilde vi komme til det legaliserede
Fællesskab (Kvindelinjen og Polyandriet), og derfra gennem
den patriarkalske Familje til det monogame Ægteskab. Denne
Klarhed og Simpelhed har vist nok øvet en forførende Magt over
de nævnte Forskere. Der kan godt være foregaaet en Udvikling,
uden at dens Bane er saa let at overskue. Og først og fremmest
bliver Spørgsmaalet, om de nævnte Stadier virkelig lade sig paavise
som historisk givne Udviklingsformer, der hver til sin Tid
have været herskende.

% --- p. 790 ---
\opage{790}Dr. Starcke nægter baade, at man fra Kvindelinjen kan
slutte tilbage til en oprindelig Promiskuitet, og at Kvindelinjen
selv har været saa almindelig, som den omtalte Teori paastaar.
Han søger dernæst at vise, at Kvindelinjen, hvor den er gældende,
kan være opstaaet af forskellige Motiver hos de forskellige Folkeslag,
og at ligeledes de andre Sæder og Skikke, man har beraabt
sig paa, maa forklares paa anden Maade end efter Teorien.

Hvad det første Punkt angaar (Kvindelinjen som Bevis paa
oprindelig Promiskuitet), kunde Forf. have anført, at der er
Forskere, som antage Kvindelinjens Almindelighed paa et vist
Trin og dog forkaste Antagelsen af en almindelig Promiskuitet.
% emphasis: „Post“ letterspaced (gaps 8, 7, 7 px at 300 dpi)
Saaledes \emph{Post} i hans »Afrikanische Jurisprudenz. Ethnologisch-juridische
Beiträge« (1887). Allerede dette tyder paa, at de to
Antagelser ikke staa og falde med hinanden. Men vigtigst er den
almindelige Betragtning, som Dr. Starcke opstiller og begrunder:
at man fra den Maade, hvorpaa Slægtskabet regnes, ikke kan
drage nogen Slutning til Ægteskabets Natur. Kvindelinjen viser
kun, at Hensynet til Faderen skydes til Side, men Spørgsmaalet
om Grunden dertil staar endnu aabent. At det ikke (overalt i
det mindste) kan have været den, at det var uvist, hvem der var
Faderen, kan ses af, at Kvindelinjen endnu den Dag i Dag gælder
under Forhold, hvor Faderen meget vel kan paavises. Den virkelige
Grund vil opdages, naar vi betænke, at Familjen paa primitive
Trin ikke eksisterer som selvstændigt Samfund, men kun som Del
af Klanen og af Stammen. Den Gruppe, til hvilken Barnet ved
Fødselen slutter sig og regnes, er Klanen. Barnet maatte regnes
% emphasis: „enten“ (gaps 6, 7, 7, 7 px) and „eller“ (gaps 7, 8, 7, 6 px) letterspaced
\emph{enten} til Faderens \emph{eller} til Moderens Klan, og af forskellige
Grunde regnedes det hos mange Folkeslag til Moderens. Men
derved berørtes det egentlige Familjeforhold ikke nødvendigvis.
Faderen kunde staa i det ømmeste Forhold til sine Børn, selv om
de regnedes til Moderens, ikke til hans Klan, og det samme
gjaldt for Moderens Vedkommende, hvor Mandslinjen herskede.
Slægtskabsbestemmelsen er altsaa en Klanbestemmelse, ved hvilken
man paa primitive Trin ikke tager Hensyn til Avlingsforholdet.
Antage vi, at Manden i Kraft af sin fysiske Styrke oprindelig har
været den herskende, og at Kvinder og Børn betragtedes som
hans Ejendom, vil Kvindelinjens Udvikling kunne forklares af de
hos forskellige Folkeslag virkende forskellige Forhold. Hos de
amerikanske Folkeslag (f. Eks. Kolumberne) spores en Sammenhæng
mellem Kvindelinjens Udvikling og Klanens fastere Udform%
% --- p. 791 ---
\opage{791}ning. Til Kvindens Klan bliver ogsaa de Børn, hun føder, og
som fra først af hænge saa nøje sammen med hende, regnede;
hun er sin Slægts Ejendom, og hendes Børn kunne igen ikke
skilles fra hende. Kvindens Værdi som Slave eller dog som Hjælp
ved Arbejdet maa her vel fastholdes; paa den vilde Klanen ikke
give Afkald, fordi hun traadte i kønslig Forbindelse med en Mand
af en anden Klan. Med Mænd af hendes egen Klan var det
hende paa Grund af det herskende Exogami forbudt at træde i
Forbindelse. Hos de afrikanske Folkeslag (f. Eks. Bechuanerne)
vil Motivet til Kvindelinjens Opstaaen være at søge i det herskende
% displaced letter: the t of „Hustru“ stands raised and apart („Hus t ru“); transcribed as the word printed
Polygami i Forbindelse med den Ordning, at hver Hustru bor for
sig, og at Børnene altsaa komme til at udgøre en Husstand
sammen med Moderen, i Modsætning til de andre Husstande, der
% a small speck at the baseline after „den“, much smaller than a full stop, taken as a blemish
med dem samlede sig om den patriarkalske Familjes fælles Hoved.
Den Stilling, Barnet faar i den patriarkalske Familje, vil her
væsentlig bero paa Moderens Rang og Værdighed. Sønnerne
bære, saa længe Faderen lever, deres Mødres Navn, og først
efter Faderens Død indtræder den fornemste Hustrus ældste Søn
i hans Stilling og Navn. Spiren til Kvindelinjen ligger i den Vægt
Moderens Stilling har faaet. — I øvrig ville ogsaa andre Forhold
blive af afgørende Betydning. Et agerdyrkende Folk vil have
Tendens til Kvindelinjen, fordi Datteren dér er af Værdi som
Arbejder; et kvægavlende Folk vil have Tendens til Mandslinjen,
fordi det der vil være fordelagtigt at sælge Datteren og faa Kvæg
i Stedet. Kvinden betragtes jo paa primitive Trin væsentlig som
en Ejendom, det gælder at anvende saa fordelagtig som muligt,
og denne Betragtning følges baade fra Faderens og fra Ægtemandens
Side.

Moderne Iagttagere og Forskere ere — i Følge Dr. Starcke
— for tilbøjelige til at overføre deres Følelser til de Vilde. For
os er Bevidstheden om at være Barnets fysiske Ophav, om at det
er Faderens Blod, der ruller i Barnets Aarer, en nødvendig Bestanddel
i den egentlige Faderfølelse. Men saaledes er det ikke
for den Vilde. Hans Forestillinger forbinde sig mere ved blot
ydre Sammenhæng end efter virkelig Kausalitet. Synsforestillingerne
ere de herskende hos os alle og ere aldeles bestemmende for den
% emphasis: „ser sammen“ (gaps 6, 7; 8, 6, 8, 9, 6 px) and „hører sammen“ (gaps 7, 6, 6, 7; 7, 6, 8, 8, 7 px) letterspaced
Vilde, saa at det som han \emph{ser sammen} (og kun det) ogsaa for
ham \emph{hører sammen}. Klanen bor sammen i sin særlige Del af
Landsbyen, og Børnene bo med deres Moder i hendes særlige
Hytte. Denne rumlige Samværen bliver af stor Betydning for
% --- p. 792 ---
\opage{792}Kvindelinjens Udvikling. Og den forklarer tillige det mærkelige
Fænomen, at den Vilde og Barbaren betragter sig som Fader,
% emphasis: „véd“ letterspaced (gaps 8, 5, 6 px at 300 dpi)
selv hvor han \emph{véd}, at han ikke er den fysiske Fader. Hans
Hustrus Børn ere stedse ogsaa hans egne Børn. Polyandri, Levirat
og Niyoga finde let og naturlig deres Forklaring ved denne Tankegang,
uden at man behøver at postulere en oprindelig Parring i
Flæng. Faderforholdet er hos de Vilde overvejende et Retsforhold,
ikke saa meget et Blodsforhold.

Foruden dette vigtige psykologiske Moment er der endnu
en Betragtning, som efter Dr. Starcke er af afgørende Betydning
for en rigtig Opfattelse af de Vildes Familjeinstitutioner. Den
Løshed og Ubundethed i kønslig Henseende, som findes hos de
Vilde, har naturlig været Iagttagere og Forskere paafaldende, men
har haft for stor Indflydelse paa deres Teorier. Tilfredsstillelse
af Kønsinstinktet spiller en stor Rolle i den Vildes Liv, men er
ogsaa let at opnaa, da der ingen snevre Skranker sættes for den
af de herskende Sæder. Men netop derfor bliver det kønslige
Forhold til et andet Individ ikke hos de Vilde af stor praktisk
Betydning, bliver ikke Grundlag for en Institution. Familjen som
Institution er derfor fra først af og i lange Tider bleven ordnet
efter helt andre Hensyn end selve de Naturtilskyndelser, der føre
Mand og Kvinde sammen. Først efter Haanden som Forestillingslivet
udvikler sig, og det ikke blot er den plumpe Rumsforestilling,
der behersker Bevidstheden, vil Faderforholdet væsentlig betragtes
som et Blodsforhold, og dermed vil Kyskhedsfordringen til Kvinden
i og snart ogsaa forud for Ægteskabet være given; fra Kvinden
vil den dernæst ogsaa blive udstrakt til Manden.

Dr. Starcke burde maaske have gennemført den her antydede
interessante psykologiske Betragtning noget videre. Ved
højere aandelig Udvikling og den dermed følgende finere Sans for
det personlige Livs Enhed og Sammenhæng vil det ogsaa blive
umuligt at opfatte Kønsinstinktet og dets Tilfredsstillelse som noget
der staar uden nogen Forbindelse med Personlighedens hele øvrige
Liv. Gennem mange fine Traade hænger denne Side af Livet
sammen med Livets andre Sider. Det viser sig navnlig derved,
at det fysiske Instinkts Tilfredsstillelse ikke er det udelukkende
Maal, men at denne Tilfredsstillelse kun staar som et Gode, naar
en dybere Sympati med et andet Individ, Glæde og Velbehag ved
% printer's error: „Kønsinktinktet“ (for „Kønsinstinktet“); transcribed as printed
dette derved faar sin Besegling og sit Udtryk. Kønsinktinktet ud%
% --- p. 793 ---
\opage{793}vikler sig da til den egentlige Elskovsfølelse, og den fysiske Forbindelse
til ogsaa at være en aandelig Forbindelse.

Man har ofte lagt for megen Romantik ind hos de Vilde
og i deres mere ubundne Kønsliv fundet den Frihed, som man
mener Elskovsfølelsen mangler i det civiliserede Liv. Dr. Starcke
lægger maaske til Gengæld for liden Romantik ind hos de Vilde.
I hans Bog fremtræder Forholdet mellem Mand og Kvinde paa
de primitive Udviklingstrin dels som et blot fysisk Kønsforhold
dels som et rent juridisk og økonomisk Forhold. Det vilde have
været af stor Interesse, om han havde benyttet sit store Materiale
til at undersøge, hvor vidt Elskovsfølelsen kan paavises hos de
primitive Mennesker, og hvilke Former den antager. Det vilde
derved sikkert have vist sig, at Forholdene dog ikke ere slet saa
forskellige som man efter Dr. Starckes Fremstilling kunde faa
Indtrykket af. Lige som det er en stor Misforstaaelse at tro, at
de Vilde leve et frit og ubundet Liv, eftersom de ere endnu større
Slaver af Traditioner, Ceremonier og Regler end de fleste civiliserede
Folk, saaledes vilde det ogsaa være en Misforstaaelse at
antage, at selv det laveste menneskelige Livstrin, vi nu kende,
aldeles ikke frembyder Følelsesfænomener, der ere beslægtede med
hvad vi nu (til Forskel for det blotte Kønsinstinkt) kalde Elskovsfølelse.
Men det vilde sikkert vise sig, at Sammenhængen mellem
Følelse og Instinkt er endnu nøjere hos de Vilde end den behøver
at være hvor en højere aandelig Udvikling har fundet Sted. Den
Vilde kender ingen Lyrik uden som Opfordring til Handling eller
som Middel til at naa sit Maal. En Forfatter, som har beskæftiget
% emphasis: „en“ letterspaced (gap 6 px at 300 dpi; the solid „en“ elsewhere on the page 0--1 px)
sig med \emph{en} betydningsfuld Side af de Vildes Aandsliv, siger:
»Hvis du beder en Indianer om en Elskovssang, vil han svare,
at en Elskovsdrik virker bedre. Den Vilde er overordentlig praktisk.
Hans Kunster, Mimik og Tegning, eksistere ikke for Kunstens egen
Skyld, men som Midler til et bestemt Formaal, som Metoder, ved
hvilke han kan opnaa noget, han ønsker. Den unge Vilde, som
Kohl [en Rejsende i Nordamerika] kendte, satte sin Lid til, at
han havde et Billede af sig selv og et Billede af sin elskede. Det
kvindelige Billedes Hjærte besmurte han med magisk Pulver; dette
var, sagde han, en almindelig Skik, og man føjede dertil elegiske
eller ondskabsfulde Sange og helt forbryderiske Besværgelser.«
% emphasis: „Andrew Lang“ letterspaced (gaps 7, 7, 11, 6, 6; 6, 7, 11 px at 300 dpi)
(\emph{Andrew Lang}: Myth, Ritual and Religion. London 1887).

Den grelle Modsætning, som de Vildes Liv alligevel uimodsigelig
udviser mellem det i mere eller mindre flygtig Opblussen
% --- p. 794 ---
\opage{794}fremtrædende Kønsinstinkt og de faste varige Former for Slægtens
Liv og Udvikling, vil kun kunne overvindes og forsones gennem
Elskovsfølelsens sunde Udvikling. Det vidunderlige ved denne
Følelse er dens Ævne til at forbinde saa mange af Livets Kræfter
om et Centrum, at forene det mest ideelle og det mest elementære
i den menneskelige Natur, og tillige dens Ævne til at undergaa
Metamorfoser gennem skiftende Skæbner og Livsperioder. Uden
denne Rigdom og denne Forvandlingsævne vilde den ikke fortjene
den fremragende Plads, som Digtekunsten har givet den.

Vi ere endnu ikke komne langt ud over de Vildes Standpunkt
hvad denne Side af Livet angaar. Det viser sig allerede
derved, at saa mange — baade af den »frie Kærligheds« Tilhængere
og af dens Modstandere — ved Ordet »Ægteskab« først og
fremmest tænke paa en i juridisk eller kirkelig Form indgaaet
Forbindelse mellem Mand og Kvinde. Dette har gjort megen bitter
Strid til en Kamp om blotte Ord og har bortledet Opmærksomheden
fra det egentlige Hovedpunkt, nemlig Trofasthedens og
Inderlighedens Mulighed og deres Betydning i Forholdet mellem
Mand og Kvinde. —

% emphasis: „Stilling“ letterspaced (gaps 8, 6, 9, 7, 9, 8, 7 px at 300 dpi)
Efter Dr. Starckes Opfattelse er der en stor og afgørende
Forskel mellem Familjens \emph{Stilling} paa primitive og paa højere
Udviklingstrin. Han viser paa en slaaende Maade, at den egentlige
Familje (Gruppen af Mand, Hustru og Børn) først ved højere
Civilisation bliver et selvstændigt, af Klaner og Stammer uafhængigt
Samfund. Hos de Vilde eksisterer Familjen kun som Del
af Klanen. Dette maatte være saaledes paa Grund af den uafbrudte
Kamp for Tilværelsen. Klanen bestaar, fordi Sammenslutning
er nødvendig i denne Kamp; i Familjen derimod nydes
de tilkæmpede og erhvervede Goder. Jo mere derfor Kampen for
den fysiske Tilværelse træder tilbage, des mere vil Klanen tabe i
Betydning, og Familjen udvikle sig selvstændig. I Familjelivet
bliver Mennesket delagtigt i en Nydelse af anden Art end den, der
opstaar ved Tilfredsstillelse af fysiske Instinkter, nemlig »den
sædelige Nydelse, som Livet for og med andre yder«. — Men
% emphasis: „Motiver“ letterspaced (gaps 10, 6, 8, 7, 8, 7 px at 300 dpi)
hvad de \emph{Motiver} angaar, som stifte Ægteskabet, er der efter Dr.
Starcke ikke og bør der (mener han) heller ikke være nogen afgørende
Forskel mellem det primitive og det egentlig humane Liv. »Vi
have set«, siger han henimod Slutningen af sin Bog, »at det oprindelig
ikke er nogen øm Følelse, i det mindste ikke Kærlighed,
som indgyder Manden Ønsket om at gifte sig, at det primitive
% --- p. 795 ---
\opage{795}Ægteskab, haardt og mørkt som det primitive Liv, udspringer af
% broken letter: the first n of „endnu“ is damaged (reads „endmu“); transcribed as the word printed
de mest konkrete og nøgterne Fornødenheder. Og endnu stedse
er det Ægteskab slet, i hvilket det erotiske er det fremtrædende
i Forholdet mellem Ægtefolkene. Det fælles Hjem, i hvilket
enhver af dem har sin Opgave, og den fælles Interesse for at faa
Børn og opdrage dem, det var de to Grundvolde, paa hvilke
Ægteskabet oprindelig blev bygget. Og af den Sympati, som
uundgaaelig udspringer af de fælles Interesser, fremvokser igen den
Kærlighed, der sluttelig fuldender og befæster Ægteskabet.« Men
vil ikke netop Elskovsfølelsen mere og mere blive en nødvendig
Forudsætning for, at en Mand og en Kvinde skulle blive til Sinds
at bygge paa hine to Grundvolde? Og kan den undværes som
et Cement, der sammenholder den Bygning, som rejses? At den
for at kunne virke saaledes maa undergaa en hel Udviklingsproces
i Forhold til de foreliggende Livsopgaver og Livsperioder, er allerede
antydet. —

Skøndt jeg har savnet noget i mine Øjne betydningsfuldt i
% printer's error: full stop after „Bog“ before lower-case „skilles“ (for a comma); transcribed as printed
Dr. Starckes Bog. skilles jeg dog fra den med fuld Anerkendelse
af den Grundighed og Skarpsindighed, hvorom den vidner. Kun
ufuldkomment har jeg i denne Anmældelse kunnet antyde den
Rigdom af Enkeltundersøgelser, som Bogen indeholder. Den er
udkommen paa Tysk, fordi en saa speciel Undersøgelse ikke kunde
gøre Regning paa tilstrækkelig mange Læsere i vort Modersmaal.
Men denne Omstændighed bør ikke udelukke Forfatteren fra tilbørlig
Anerkendelse fra hans Landsmænd.

% p. 795: signature set right in small caps; below it a clearly printed centred double rule (two thin lines c. 5 px apart at 300 dpi, c. 2.6 cm long, 305 px), the end-of-article rule; the rest of the page is blank; nothing of the next article („Han bare laa og smilte“, p. 796) on this page
\begin{flushright}\textsc{Harald Høffding.}\end{flushright}
\begin{center}\rule{2.6cm}{0.4pt}\\[-10pt]\rule{2.6cm}{0.4pt}\end{center}

\end{document}
